\chapter*{Notation and conventions}
\addcontentsline{toc}{chapter}{Notation and conventions}

This chapter fixes one set of conventions for the whole of the notes. They are the ones most common in the string-compactification literature, which are not always those used on the board; where a lecture writes a formula in another convention, it is converted to the one below.

\paragraph{Units and approximate equalities.} We use natural units, $\hbar = c = 1$, so that masses, energies and inverse lengths are measured in $\mathrm{GeV}$. An equality sign is exact. The sign ``$\simeq$'' means equality up to numerical factors of order one (such as $2$, $\pi$ or $\sqrt{4\pi}$) that are dropped when only the parametric dependence on $g_s$ and $\Vol$ matters, and ``$\sim$'' means ``of the order of''. When a relation is first derived its factors are kept, and it is then quoted with ``$\simeq$''.

\paragraph{Metric and curvature.} Metrics have mostly-plus signature, $\eta_{MN} = \mathrm{diag}(-1,+1,\dots,+1)$, as in Polchinski~\cite{Polchinski}, Tong~\cite{Tong}, Wess and Bagger, and most of the compactification literature. The Riemann tensor is $R^\rho{}_{\sigma\mu\nu} = \partial_\mu\Gamma^\rho_{\nu\sigma} - \partial_\nu\Gamma^\rho_{\mu\sigma} + \Gamma^\rho_{\mu\lambda}\Gamma^\lambda_{\nu\sigma} - \Gamma^\rho_{\nu\lambda}\Gamma^\lambda_{\mu\sigma}$, the Ricci tensor is $R_{\mu\nu} = R^\rho{}_{\mu\rho\nu}$, and a sphere has positive Ricci scalar. With these choices the four-dimensional Einstein--Hilbert action is
\begin{equation*}
  S_{\mathrm{EH}} = \frac{\Mp^2}{2}\int\dd^4x\,\sqrt{-g}\,R .
\end{equation*}
For a Lorentzian metric we write $\sqrt{-g}$, for the Riemannian metric of the compact space $\sqrt{g}$.

\paragraph{Spaces, coordinates and indices.} Ten-dimensional spacetime is $X_{10} = \R^{1,3}\times Y_6$, with $Y_6$ a compact six-dimensional manifold and coordinates $X^M = (x^\mu, y^m)$.
\begin{itemize}
  \item Capital Latin indices $M,N,P = 0,\dots,9$ are ten-dimensional.
  \item Greek indices $\mu,\nu,\rho = 0,\dots,3$ are four-dimensional.
  \item Lower-case Latin indices $m,n,p = 4,\dots,9$ are real indices on $Y_6$.
  \item The indices $i,j = 1,2,3$ and $\bar\imath,\bar\jmath$ label holomorphic and antiholomorphic coordinates $z^i$, $\bar z^{\bar\jmath}$ on $Y_6$ when it is a complex manifold.
\end{itemize}
The ten-dimensional metric is $G_{MN}$, with determinant $G$ and Ricci scalar $R_{10}$. On a product space it splits into the four-dimensional metric $g_{\mu\nu}$ (determinant $g_4$, Ricci scalar $R_4$) and the internal metric $g_{mn}$ (determinant $g_6$, Ricci scalar $R_6$). The volume of the compact space is $\mathrm{Vol}(Y_6) = \int_{Y_6}\dd^6y\,\sqrt{g_6}$. The holonomy group of $Y_6$ is $\Hol(Y_6)$, and $\one$ is the identity.

\paragraph{String length and string scale.} The only dimensionful parameter of string theory is $\alpha'$, related to the string tension by $T = 1/(2\pi\alpha')$. Following the lecture and the moduli-stabilisation literature we define
\begin{equation*}
  \ls = 2\pi\sqrt{\alpha'} ,\qquad \Ms = \frac{1}{\ls} .
\end{equation*}
Polchinski and Tong use $\ls = \sqrt{\alpha'}$ instead, which is smaller by $2\pi$. The virtue of the factor $2\pi$ is that it absorbs the powers of $2\pi$ of the ten-dimensional gravitational coupling, $2\kappa_{10}^2 = (2\pi)^7\alpha'^4 = \ls^8/2\pi$, and of the D$p$-brane tensions, $T_p = 2\pi\Ms^{p+1}/g_s$. The price is that T-duality, which exchanges a circle of radius $R$ with one of radius $\alpha'/R$, reads $R\to\alpha'/R = 1/(4\pi^2\Ms^2R)$ (\cref{eq:L2-Tduality}); the self-dual radius is $\sqrt{\alpha'} = \ls/2\pi$.

\paragraph{Dilaton, string coupling and frames.} The dilaton $\phi$ is a field. Its vacuum expectation value is $\phi_0 \equiv \langle\phi\rangle$, and the string coupling is
\begin{equation*}
  g_s = e^{\phi_0} .
\end{equation*}
The \emph{string frame} metric $G_{MN}$ is the one to which the string couples on the worldsheet; in it the gravitational action carries a factor $e^{-2\phi}$ in front of the Ricci scalar (\cref{eq:L1-S10}). The ten-dimensional \emph{Einstein frame} metric is
\begin{equation*}
  G^E_{MN} = e^{-(\phi-\phi_0)/2}\,G_{MN} ,
\end{equation*}
in which the gravitational action takes the Einstein--Hilbert form $\frac{1}{2\kappa_{10}^2 g_s^2}\int\dd^{10}X\sqrt{-G_E}\,R_{10}[G_E]$. Because $\phi_0$ is subtracted, the two frames coincide in a vacuum with constant dilaton, $\phi = \phi_0$. In particular a volume measured in string units is the same number in both frames, and the relations between scales derived in \cref{ch:L1}, $\Mp^2 = 4\pi\Vol\Ms^2/g_s^2$ and $\MKK \simeq \Ms/\Vol^{1/6}$, hold in either.

Much of the moduli-stabilisation literature uses instead $G^E_{MN} = e^{-\phi/2}G_{MN}$, without the subtraction. In a vacuum the two Einstein frames then differ by the constant factor $g_s^{-1/2}$, so that a dimensionless volume in that convention is $\Vol' = g_s^{-3/2}\,\Vol$, and the same relations read $\Ms = g_s^{1/4}\Mp/\sqrt{4\pi\Vol'}$ and $\MKK\simeq\Mp/(\sqrt{4\pi}\,\Vol'^{2/3})$. Formulas taken from such references are converted with $\Vol' = g_s^{-3/2}\Vol$.

\paragraph{Volume.} $\Vol$ is the dimensionless volume of $Y_6$ in units of $\ls^6$, measured with the Einstein-frame metric:
\begin{equation*}
  \mathrm{Vol}(Y_6) = \Vol\,\ls^6 .
\end{equation*}

\paragraph{Planck mass.} $\Mp$ is the reduced Planck mass,
\begin{equation*}
  \Mp = (8\pi G_N)^{-1/2} \simeq 2.4\times10^{18}\GeV ,
\end{equation*}
the coefficient of the four-dimensional Einstein--Hilbert action above. Four-dimensional actions are written in the four-dimensional Einstein frame, in which this coefficient is constant.

\paragraph{Counting supersymmetries.} $\mathcal N$ counts supersymmetry generators in units of the smallest spinor of the spacetime dimension. In ten dimensions this is a Majorana--Weyl spinor, with $16$ real components, so type IIA and IIB supergravity, with two of them, have ten-dimensional $\mathcal N=2$ and $32$ real supercharges. The two have opposite chirality in type IIA, and the same chirality in type IIB, which is therefore chiral. In four dimensions the unit is a Majorana spinor, equivalently a Weyl spinor and its conjugate, with $4$ real components, so that four-dimensional $\mathcal N = 8,4,2,1$ correspond to $32,16,8,4$ real supercharges. A subscript $\pm$ on a ten-dimensional spinor, as in $\psi_{M+}$ or $\lambda_-$, gives its chirality, the eigenvalue $\pm1$ of $\Gamma_{11} = \Gamma^0\Gamma^1\cdots\Gamma^9$, which squares to one in our signature.

\paragraph{Forms.} A $p$-form $C_p = \frac{1}{p!}C_{M_1\dots M_p}\,\dd X^{M_1}\wedge\dots\wedge\dd X^{M_p}$ is identified with its totally antisymmetric components, e.g.\ $C_2\leftrightarrow C_{MN} = C_{[MN]}$. The NS--NS two-form is $B$, with field strength $H = \dd B$; the R--R forms are $C_p$, with field strengths $F_{p+1} = \dd C_p + \dots$. A form $\omega$ on $Y_6$ is harmonic if $\Delta\omega = 0$, with $\Delta = \dd\dd^\dagger + \dd^\dagger\dd$ the Laplacian of the internal metric. On a complex manifold a $(p,q)$-form has $p$ holomorphic and $q$ antiholomorphic indices.

\paragraph{Moduli.} The Hodge numbers $h^{p,q}$ of $Y_6$ count its linearly independent harmonic $(p,q)$-forms; for a Calabi--Yau threefold the non-trivial ones are $h^{1,1}$ and $h^{2,1}$ ($=h^{1,2}$). A generic modulus is $\phi^I$; in discussing the couplings of moduli to matter, $\phi$ without index also stands for a generic canonically normalised modulus. The Kähler moduli are $T_i$, $i = 1,\dots,h^{1,1}$, and the complex-structure moduli $U_a$, $a = 1,\dots,h^{2,1}$. The index $i$ of $T_i$ is a label, not a holomorphic index; the two are never needed in the same formula, and the context tells which is meant. The axio-dilaton is
\begin{equation*}
  S = e^{-\phi} + i\,C_0 ,\qquad \langle\mathrm{Re}\,S\rangle = \frac{1}{g_s} ,
\end{equation*}
as in the lecture and in the Large Volume Scenario literature. Much of the flux-compactification literature (GKP, KKLT, F-theory) uses instead $\tau = C_0 + i\,e^{-\phi} = i\bar S$, on which the $SL(2,\Z)$ S-duality of type IIB acts by fractional linear transformations; formulas taken from there are converted with $\tau = i\bar S$. With an index, $\tau_i$ will denote the volumes of four-cycles, the real parts of the $T_i$. Gauge couplings are $\alpha = g^2/4\pi$; the hypercharge coupling is GUT-normalised, $\alpha_1 = \frac53\alpha_Y$.

\paragraph{Supergravity.} Four-dimensional $\mathcal N=1$ supergravity, with Kähler potential $K$ and superpotential $W$, is written in the conventions of Wess and Bagger.
